Can a marketing-specialized company memory improve next-action recommendations, or is strong generic context sufficient? We separate the policy value recoverable from retrieved evidence from the grounded feasible value actually returned by a fixed reader. A frozen benchmark contains 120 fictional company worlds, 360 multi-account cases, three compound stress levels, six executed context treatments and 4,320 actual GPT-6.1 Sol/Claude Opus 5.5 responses. On positive cases, general company context yields 47.9\% of the full-information optimum for both readers; GTM context yields 88.4\% for GPT and 88.2\% for Opus. Paired gains are 40.6 percentage points (95\% company-bootstrap interval 36.7--44.4) and 40.3 points (36.3--44.3). No company evidence yields zero grounded positive value, as expected under the evidence-required contract. The advantage is principally information selection: a generic query-word selector produces exactly the same GTM packets, and wider generic contexts close the retrieval gap. Full current context performs better, at greater API cost. On a separate one-company raw-document/wiki development set, the general wiki answers 22/22 questions correctly versus 21/22 for the marketing wiki. The results support task-aware evidence selection under a budget, not universal vertical-memory superiority, a proprietary architecture requirement or measured revenue improvement. Data, code, raw attempts, uncertainty analysis and publication figures are packaged for reproduction.
